\chapter{Theory}
\label{ch:architecture}

\section{Optical Principles and Numerical Simulation}
在這個章節中，我會詳細推導並說明光在矽基整合光學波導中的傳播物理，以及模態分波多工器運作的理論基礎，
最後會介紹數值模擬的基礎理論。我們將從電磁場數學模型出發，分析奈米級波導的邊界條件與模態行為，
作為後續反向設計演算法與模型預測的依據。
\subsection{Waveguide Physics and Eigenmode Theory}
首先，在波導中，光會被侷限在折射率較高的矽核心中，並透過全反射(Total internal reflection, TIR)進行傳播。
關於介電質波導中電磁場行為與模態特徵值問題的完整理論架構，本節主要參考 A. Yariv 的經典著作 [29] 所建立的物理模型。
在無自由電荷($\rho = 0$)且無源($\vec{J} = 0$)的介電質波導環境中，
我們可以用馬克士威方程(Maxwell equations)中的法拉第電磁感應定律(3.1)與馬克士威-安培定律(3.2)微分式
來描述光的傳播行為，如下所示：
\begin{align}
\nabla \times \vec{E} &= -\frac{\partial \vec{B}}{\partial t}
\label{eq:maxwell1} \\
\nabla \times \vec{H} &= \frac{\partial \vec{D}}{\partial t}
\label{eq:maxwell2}
\end{align}
這時候引入了材料的介電質($\epsilon$)和導磁率($\mu$)，假設光場在具有時諧性(Time-Harmonic)，則(3.1)和(3.2)
馬克士威方程可以寫成頻域複數形式:
\begin{align}
\nabla \times \vec{E} &= -i\omega\mu\vec{H}
\label{eq:freq_maxwell1} \\
\nabla \times \vec{H} &= i\omega\epsilon\vec{E}
\label{eq:freq_maxwell2}
\end{align}
接著，我們再對馬克士威方程頻域複數形式(3.3)(3.4)取旋度，再均勻的介質中，由於$\nabla\cdot\vec{E} = 0$ ，
因此，我們即可推導出經典的亥姆霍茲方程(Helmholtz Equation):
\begin{align}
\nabla^2\vec{E} + k_0^2 n^2(\vec{r})\vec{E} = 0
\label{eq:helmholtz}
\end{align}
其中，$k_0 = \frac{\omega}{c}$為真空中的波數，介電常數與折射率($n$)關係為$\epsilon = n^2\epsilon_0$，
因此光子能夠被侷限在波導內部傳播，就是因為波導核心層與包層之間的折射率差而產生的全反射(TIR)，在本篇矽基波導的核心
材料為矽(Silicon, Si)，其折射率($n_{\mathrm{Si}}$)約為3.48; 而包層為二氧化矽(Silicon dioxide, $\mathrm{SiO_2}$)，
其折射率($n_{\text{SiO}_2}$)約為1.44。

當波導結構在傳播方向($z$軸)上具備平移對稱性(Translation symmetry)時，我們可以透過將電場進行解分離變數，
寫成沿著$z$方向傳播的形式:
\begin{align}
\vec{E} = \vec{E}(x, y)e^{-i\beta z}
\label{eq:Translation symmetry in Z}
\end{align}

接著將$z$方向傳播形式(3.6)帶入亥姆霍茲方程(3.5)後，對$z$做二次偏微分會得到$-\beta^2$(也就是$\frac{\partial^2}{\partial z^2} \rightarrow -\beta^2$)，
整理後我們可以得到橫截面($xy$平面)上的二維波動方程:
\begin{align}
\left( \frac{\partial^2}{\partial x^2} + \frac{\partial^2}{\partial y^2} \right)\vec{E} + k_0^2 n^2(\vec{r})\vec{E} = \beta^2\vec{E}
\label{eq:xy 2D Wave Equation}
\end{align}

這邊可以定義成一個線性算子(Linear operator)$L$:
\begin{align}
L \triangleq \left( \frac{\partial^2}{\partial x^2} + \frac{\partial^2}{\partial y^2} \right) + k_0^2 n^2(\vec{r})
\label{eq:linear_operator}
\end{align}

因此可以將整個光學波導的模態傳播問題簡化為標準的特徵值問題(Eigenvalue Problem, EVP):
\begin{align}
L\vec{E} = \beta^2\vec{E}
\label{eq:WG Eigenvalue Problem}
\end{align}

在這個式子中，模態(Modes)的空間場分佈就會等於該算子 $L$ 的特徵函數(Eigenfunction)，
而傳播常數的平方 $\beta^2$ 則會等於其對應的特徵值(Eigenvalue)。這個光學特徵值問題在數學形式上，
與量子力學中的一維有限位能井問題(Particle in a box problem)有著異曲同工之妙，都必須在邊界條件下求得特定波數與對應的
空間分布，一維定態薛丁格方程(Schrödinger equation)表示如下：
\begin{align}
-\frac{\hbar^2}{2m}\nabla^2\psi + V(\vec{r})\psi = E\psi
\label{eq:1D schrodinger equation}
\end{align}
此式的波函數($\psi$)、位能分布(V($\vec{r}$))與能量特徵值($E$)，分別對應到了波導方程(3.7)中的電場($\vec{E}$)、折射率分布($k_0^2 n^2(\vec{r})$)
以及傳播常數的平方($\beta^2$)。在數學形式的對等性中，我們可以得出光波導兩個核心物理特質的結論:模態量子化(Quantization)與能量侷限性
(Confinement)，只有特定且不連續的$\beta$值才能夠穩定存在於波導中，且折射率的對比愈大，位能井就愈深，則波導對光場的侷限性就越好。

另外，每個模態的傳播特徵可以由有效折射率(Effective Refractive Index)來定義:
\begin{align}
n_{eff} \triangleq \frac{\beta}{k_0} 
\label{eq:effective refractive index}
\end{align}
因此，對於能夠在波導中穩定傳播的束縛模態，其有效折射率($n_{eff}$)必須介於核心折射率($n_2$)和包層折射率($n_1$)之間，也就是，
\begin{align}
n_{1} < n_{eff} < n_{2}
\label{eq:effective refractive index range}
\end{align}
其原因可由電磁場在核心層與包層中的分布特性來理解。在核心層內，光場必須呈現簡諧波動解(Sinusoidal solution)傳播，因此橫向波數需滿足
\begin{align}
h^2\triangleq k_0^2 n_2^2 - \beta^2 > 0
\label{eq:Sinusoidal solution}
\end{align}
由此可得$\beta < n_2 k_0$，也就是$n_{eff} < n_2$。\
另一方面，在包層中光場需以指數衰減的形式向外傳播，也就是漸逝場(Evanescent Field)的形式向外衰減，以確保光能量被侷限於波導內，
因此衰減係數必須滿足
\begin{align}
q^2 = \beta^2 - k_0^2 n_1^2 > 0
\label{eq:Evanescent Field solution}
\end{align}
由此可得$\beta > n_1 k_0$，也就是$n_{eff} > n_1$。
因此，束縛模態的有效折射率必須介於核心層與包層折射率之間，即使是結構較為複雜的反向設計波導，仍然需要滿足以上條件。\

接著，本節參考 A. Yariv 經典光電物理模型中所建立的分析架構 [29]，以一維平板波導(Slab Waveguide)結構為例，
考慮只有具備$y$軸分量的橫電模態(Transverse Electric Modes, TE modes)，
其電場可以由以下公式表示:
\begin{align}
E_y(x,z,t)=E_m(x)e^{i(\omega t-\beta z)}
\label{eq:TE mode electric field in slab waveguide}
\end{align}
代入二維波動方程(3.7)可以簡化為一維常微分方程，如下式所示:
\begin{align}
\frac{d^{2}}{d x^{2}}E_{m} + \left[ k_{0}^{2}n^{2}(x) - \beta^{2} \right] E_{m} = 0
\label{eq:TE mode 1D differential equation in slab waveguide}
\end{align}
其中，核心層與包層之折射率分布表示如下:
\begin{align}
n(x) = \begin{cases} n_1, & |x| \ge \frac{d}{2} \\ n_2, & |x| < \frac{d}{2} \end{cases}
\end{align}
由於平板波導的幾何結構在中心軸($x$=0)具有對稱性，因此特徵解可以解為偶對稱(Symmetric)與奇對稱(Anti-Symmetric)的解，
其電場特徵函數 $E_m(x)$ 表示如下：
\begin{itemize}
    \item 偶對稱模態(如TE$_0$, TE$_2$):
    \begin{align}
E_{m}(x) = \begin{cases} C \cos(hx), & |x| < \frac{d}{2} \\ B e^{-q |x|}, & |x| \ge \frac{d}{2} \end{cases}
\label{eq:TE mode symmetric solution}
\end{align}
    \item 奇對稱模態(如TE$_1$, TE$_3$):
    \begin{align}
E_{m}(x) = \begin{cases} C \sin(hx) \cdot \text{sgn}(x), & |x| < \frac{d}{2} \\ B e^{-q |x|} \cdot \text{sgn}(x), & |x| \ge \frac{d}{2} \end{cases}
\label{eq:TE mode anti symmetric solution}
\end{align}
\end{itemize}
此處橫向波數 $h^2 \triangleq k_0^2 n_2^2 - \beta^2 > 0$ 且包層消散係數 $q^2 \triangleq \beta^2 - k_0^2 n_1^2 > 0$ 。\\
根據電磁學邊界條件，在介電質交界面（$x = \pm d/2$）上，切線電場分量 $E_y$ 與切線磁場分量 $H_z$ 皆必須滿足連續條件。
透過法拉第電磁感應定律可導出磁場分量關係式 $H_{z} = \frac{1}{-i\omega\mu}\frac{\partial E_{y}}{\partial x}$ 。
將兩者的電場解代入並應用 $x = d/2$ 處的連續條件，可分別得到對應的齊次線性方程組：
\begin{itemize}
    \item 偶對稱邊界條件:
    \begin{align}
    B e^{-q\frac{d}{2}} - C \cos\left(h\frac{d}{2}\right) &= 0 \\
    B q e^{-q\frac{d}{2}} - C h \sin\left(h\frac{d}{2}\right) &= 0
    \end{align}
    \item 奇對稱邊界條件:
    \begin{align}
    B e^{-q\frac{d}{2}} - C \sin\left(h\frac{d}{2}\right) &= 0 \\
    B q e^{-q\frac{d}{2}} + C h \cos\left(h\frac{d}{2}\right) &= 0
    \end{align}
\end{itemize}

若要使上述方程組存在非平凡解（Nontrivial solutions），其係數矩陣的行列式值必須嚴格等於零 。
展開行列式並消去共同項 $e^{-qd/2}$ 後，即可成功導出兩種模態的特徵方程式（Characteristic Equations） ：
\begin{align}
    \begin{aligned}
    \text{Symmetric (偶對稱):} \quad & q = h \tan\left(\frac{hd}{2}\right) \\
    \text{Anti-symmetric (奇對稱):} \quad & q = -h \cot\left(\frac{hd}{2}\right)
    \end{aligned}
\end{align}

將 $h$ 與 $q$ 的定義代回上式(3.24)，即可得到最終用來求解傳播常數 $\beta$ 的物理關係式 ：
\begin{align}
    \begin{aligned}
    \text{Symmetric:} \quad & \sqrt{\beta^{2}-{k_{0}}^{2}{n_{1}}^{2}} = \sqrt{k_{0}^{2}n_{2}^{2}-\beta^{2}} \tan\left[\sqrt{k_{0}^{2}n_{2}^{2}-\beta^{2}}\frac{d}{2}\right] \\
    \text{Anti-symmetric:} \quad & \sqrt{\beta^{2}-{k_{0}}^{2}{n_{1}}^{2}} = -\sqrt{k_{0}^{2}n_{2}^{2}-\beta^{2}} \cot\left[\sqrt{k_{0}^{2}n_{2}^{2}-\beta^{2}}\frac{d}{2}\right]
    \end{aligned}
\end{align}

在實際的光電工程分析中，此方程式可以藉由數值計算工具來進行求解，來計算各階模態精確解的特徵值$\beta$，
以此來描繪出如$\text{TE}_0, \text{TE}_1, \text{TE}_2, \text{TE}_3$ 等完整模態橫向波函數分佈。


\subsection{Waveguide Discontinuities and Device Principles}
\subsection{Finite-Difference Time-Domain (FDTD) Method}